% ============================================================
%  Section 3: How to Create a Data Warehouse
% ============================================================
\section{Creating a Data Warehouse}

% ------------------------------------------------------------
\begin{frame}
  \frametitle{Overview: The DWH Development Process}

  Building a DWH is an iterative, multi-phase project.

  % PICTURE PLACEHOLDER
  \begin{alertblock}{Picture: DWH Development Lifecycle}
    \textit{Add a circular/iterative lifecycle diagram with phases:
    (1) Requirements Analysis $\rightarrow$ (2) Source Analysis
    $\rightarrow$ (3) Data Modelling $\rightarrow$ (4) ETL Design \& Dev
    $\rightarrow$ (5) Testing $\rightarrow$ (6) Deployment
    $\rightarrow$ (7) Monitoring \& Maintenance $\rightarrow$ back to (1).}
  \end{alertblock}

  \begin{block}{Key Principle}
    Do \textbf{not} try to build everything at once.
    Start with a focused scope (one subject area / one data mart)
    and expand iteratively.
  \end{block}

\end{frame}

% ------------------------------------------------------------
\begin{frame}
  \frametitle{Phase 1: Requirements Analysis}

  \begin{block}{Goal}
    Understand \textbf{what questions} the business needs to answer.
  \end{block}

  \begin{itemize}
    \item Interview stakeholders: managers, analysts, power users
    \item Identify \textbf{Key Performance Indicators} (KPIs)
    \item Determine the required \textbf{granularity} of data
    \item Define the \textbf{history depth} needed (1 year? 10 years?)
    \item Understand refresh frequency: daily batch? Near real-time?
  \end{itemize}

  \begin{examples}{Questions to ask stakeholders}
    \begin{itemize}
      \item \textit{``What decisions does your department make weekly?''}
      \item \textit{``What data do you currently look up manually?''}
      \item \textit{``What would you do if you had this data instantly?''}
    \end{itemize}
  \end{examples}

\end{frame}

% ------------------------------------------------------------
\begin{frame}
  \frametitle{Phase 2: Source System Analysis}

  \begin{block}{Goal}
    Understand and document the \textbf{data landscape}: what data exists,
    where it lives, and how good it is.
  \end{block}

  \begin{itemize}
    \item Inventory all source systems and their owners
    \item Understand the data model of each source (schema, relationships)
    \item Assess \textbf{data quality}: completeness, accuracy, consistency
    \item Identify \textbf{change detection}: how to know what changed since last load?
      (timestamps, CDC, log-based, full-diff)
    \item Negotiate \textbf{data access}: APIs, DB connections, file exports
  \end{itemize}

  \begin{alertblock}{Common Surprise}
    Data quality in source systems is \textbf{almost always worse}
    than stakeholders expect. Plan for data cleansing effort.
  \end{alertblock}

\end{frame}

% ------------------------------------------------------------
\begin{frame}
  \frametitle{Kimball's Dimensional Modelling Steps}

  Ralph Kimball defined a 4-step process for designing the dimensional model:

  \begin{enumerate}
    \item \textbf{Select the business process}\\
      e.g.\ ``order management'', ``customer payments''
    \item \textbf{Declare the grain}\\
      The most atomic level of data to store.\\
      e.g.\ ``one row per order line item''
    \item \textbf{Identify the dimensions}\\
      Who, what, when, where, why -- the context of each fact.
    \item \textbf{Identify the facts}\\
      The numeric measurements: revenue, quantity, duration, \ldots
  \end{enumerate}

  \begin{alertblock}{Grain First!}
    Declaring the grain \textbf{before} choosing facts or dimensions
    is the most important step. Everything else follows from it.
  \end{alertblock}

\end{frame}

% ------------------------------------------------------------
\begin{frame}
  \frametitle{Phase 3: Data Modelling -- Star Schema Design}

  \begin{block}{Designing the Fact Table}
    \begin{itemize}
      \item One row = one \textbf{event at the declared grain}
      \item Contains foreign keys to all dimension tables
      \item Contains the \textbf{additive measures} (facts): amount, count, duration
      \item Avoid storing textual descriptions in the fact table
    \end{itemize}
  \end{block}

  \begin{block}{Additive vs.\ Non-Additive Facts}
    \begin{tabularx}{\linewidth}{@{}l X@{}}
      \toprule
      \textbf{Type}      & \textbf{Description} \\
      \midrule
      Additive           & Can SUM across all dimensions (e.g.\ revenue) \\
      Semi-additive      & Can SUM across some dimensions (e.g.\ balance) \\
      Non-additive       & Cannot SUM meaningfully (e.g.\ ratio, price) \\
      \bottomrule
    \end{tabularx}
  \end{block}

\end{frame}

% ------------------------------------------------------------
\begin{frame}
  \frametitle{Phase 4: ETL Development}

  \begin{block}{Extract}
    \begin{itemize}
      \item Connect to source system and pull data (full or incremental)
      \item \textbf{Change Data Capture (CDC)}: only extract what changed
      \item Common CDC methods: timestamp columns, database triggers, log mining
    \end{itemize}
  \end{block}

  \begin{block}{Transform}
    \begin{itemize}
      \item Cleanse: handle nulls, remove duplicates, fix encodings
      \item Harmonise: map codes to uniform values (e.g.\ ``M'' and ``male'' $\to$ ``Male'')
      \item Derive new attributes (e.g.\ age from birth date)
      \item Apply \textbf{surrogate keys}: replace natural keys with integer DWH keys
    \end{itemize}
  \end{block}

  \begin{block}{Load}
    \begin{itemize}
      \item Insert new dimension records; apply SCD logic for changes
      \item Insert new fact rows (usually append-only)
    \end{itemize}
  \end{block}

\end{frame}

% ------------------------------------------------------------
\begin{frame}
  \frametitle{Surrogate Keys}

  \begin{block}{Natural Key vs.\ Surrogate Key}
    \begin{itemize}
      \item \textbf{Natural key}: the business identifier from the source system
        (e.g.\ customer number, product code)
      \item \textbf{Surrogate key}: a meaningless integer generated by the DWH
        (e.g.\ 1, 2, 3, \ldots)
    \end{itemize}
  \end{block}

  \begin{block}{Why Use Surrogate Keys?}
    \begin{itemize}
      \item Natural keys can change in source systems
      \item Surrogate keys remain stable -- essential for SCD Type 2
      \item Smaller keys = faster joins (integer vs.\ string)
      \item Decouples the DWH from source system changes
    \end{itemize}
  \end{block}

\end{frame}

% ------------------------------------------------------------
\begin{frame}
  \frametitle{Data Quality in the ETL Process}

  \begin{block}{Common Data Quality Issues}
    \begin{itemize}
      \item \textbf{Completeness}: missing or NULL values
      \item \textbf{Consistency}: same entity represented differently
        (``GmbH'' vs.\ ``Gesellschaft mbH'')
      \item \textbf{Accuracy}: values that are wrong or outdated
      \item \textbf{Timeliness}: data that arrives too late to be useful
      \item \textbf{Duplicates}: same record loaded multiple times
    \end{itemize}
  \end{block}

  \begin{block}{Strategies}
    \begin{itemize}
      \item Validate at load time; reject or quarantine bad records
      \item Use a \textbf{``bad data'' / error table} for rejected rows
      \item Implement data quality metrics and monitoring dashboards
      \item Work with source system owners to fix issues at the source
    \end{itemize}
  \end{block}

\end{frame}

% ------------------------------------------------------------
\begin{frame}
  \frametitle{Phase 5--7: Testing, Deployment, and Monitoring}

  \begin{block}{Testing}
    \begin{itemize}
      \item \textbf{Row count checks}: did all expected rows load?
      \item \textbf{Reconciliation}: totals in DWH match totals in source
      \item \textbf{Regression tests}: check that historical data has not changed
      \item \textbf{Business validation}: have domain experts sign off on reports
    \end{itemize}
  \end{block}

  \begin{block}{Deployment \& Monitoring}
    \begin{itemize}
      \item Automate ETL jobs with a scheduler (e.g.\ Apache Airflow, cron)
      \item Monitor job run times and failure rates
      \item Alert on late arrivals, row count deviations, or quality failures
      \item Maintain \textbf{lineage documentation}: where does each column come from?
    \end{itemize}
  \end{block}

\end{frame}

% ------------------------------------------------------------
\begin{frame}
  \frametitle{Common DWH Tools (Overview)}

  \begin{block}{ETL / ELT Tools}
    \begin{itemize}
      \item Traditional: \textbf{Informatica PowerCenter}, \textbf{SSIS} (Microsoft),
        \textbf{Talend}
      \item Modern / cloud-native: \textbf{dbt} (data build tool),
        \textbf{Apache Spark}, \textbf{Fivetran}
    \end{itemize}
  \end{block}

  \begin{block}{DWH Storage Engines}
    \begin{itemize}
      \item On-premise: \textbf{Teradata}, \textbf{SAP BW}, \textbf{Oracle DW}
      \item Cloud: \textbf{Snowflake}, \textbf{Google BigQuery},
        \textbf{Amazon Redshift}, \textbf{Azure Synapse}
    \end{itemize}
  \end{block}

  \begin{alertblock}{Note (!)}
    Tool knowledge is supplementary. The concepts and architecture
    are what matter for this course.
  \end{alertblock}

\end{frame}
